\documentclass[11pt,a4paper]{article}

\usepackage[utf8]{inputenc}
\usepackage[T1]{fontenc}
\usepackage[spanish,es-tabla]{babel}
\usepackage[margin=2.3cm]{geometry}
\usepackage{booktabs}
\usepackage{array}
\usepackage{xcolor}
\usepackage{hyperref}
\usepackage{enumitem}

\definecolor{azuluct}{RGB}{20, 65, 130}
\definecolor{grisoscuro}{RGB}{50, 60, 70}

\hypersetup{
    colorlinks=true,
    linkcolor=azuluct,
    urlcolor=azuluct,
    pdftitle={Guia de Estudio y Defensa - Proyecto 2 Interfaces Graficas},
    pdfauthor={Fabian Garcia}
}

\setlength{\parskip}{0.5em}
\setlength{\parindent}{0pt}

\title{\textbf{\textcolor{azuluct}{Guía de Arquitectura, Flujos y Defensa}}\\
\Large Proyecto \#2 (Unidad II) --- DataLogger Ambiental con Lazarus, Flask, SQLite y Matplotlib\\
\normalsize Interfaces Gráficas (INFO1128) --- Universidad Católica de Temuco}
\author{Fabián García}
\date{Noviembre 2026}

\begin{document}

\maketitle
\tableofcontents
\vspace{1em}
\hrule
\vspace{1em}

\section{Visión General: ¿Cómo funciona todo el proyecto de punta a punta?}

El proyecto simula una red de \textbf{5 Estaciones de Monitoreo Ambiental (UMA 1 a UMA 5)} que miden contaminación por material particulado ($\text{PM}_{1.0}$, $\text{PM}_{2.5}$ y $\text{PM}_{10}$), conteo de partículas en 6 tamaños ($0.3\,\mu\text{m}$ a $10\,\mu\text{m}$), temperatura y humedad relativa.

El flujo completo de los datos atraviesa \textbf{4 componentes} conectados en serie:

\begin{enumerate}[leftmargin=1.5em]
    \item \textbf{Simulador Serial en Python (\texttt{2\_simulador\_serial/simulador\_serial.py}):}
    Actúa como los sensores físicos. En un ciclo continuo va rotando entre las estaciones \texttt{UMA 1} a \texttt{UMA 5}, genera una trama de texto con \textbf{12 números enteros separados por comas} y la escribe en el puerto serial \texttt{COM2} a \texttt{9600} baudios.
    
    \item \textbf{Puerto Serial Virtual en VSPE (\texttt{Connector} en \texttt{COM2}):}
    En Windows, un puerto serial físico normal solo puede ser abierto por un programa a la vez. El dispositivo tipo \textbf{\texttt{Connector}} de VSPE crea un puente virtual en \texttt{COM2} que permite que Python (transmisor) y Lazarus (receptor) estén conectados al mismo puerto \texttt{COM2} simultáneamente.
    
    \item \textbf{Interfaz Gráfica DataLogger en Lazarus Pascal (\texttt{3\_lazarus\_datalogger/main.pas}):}
    Escucha el puerto \texttt{COM2} mediante el componente \texttt{TLazSerial} (\texttt{PuertoSerial}). Cada vez que entra una línea por el puerto:
    \begin{itemize}
        \item Separa los 12 números por las comas y actualiza la memoria de esa estación UMA.
        \item Si la trama corresponde a la UMA seleccionada en los \texttt{RadioButtons}, actualiza los gráficos en tiempo real (\textbf{Scrolling Left} en el gráfico de líneas y las 6 barras en el histograma) junto con la Temperatura y Humedad.
        \item Si la casilla \textbf{\texttt{Envía a Flask}} está marcada, empaqueta los 12 datos en un objeto \textbf{JSON} y los envía mediante una petición \textbf{HTTP POST} a \texttt{http://127.0.0.1:7070/data}.
    \end{itemize}
    
    \item \textbf{Servidor Web Flask + SQLite + Matplotlib (\texttt{1\_servidor\_flask/servidor.py}):}
    Corre en el puerto \texttt{7070}. Cumple tres tareas centrales:
    \begin{itemize}
        \item \textbf{Ingesta y Base de Datos SQLite (\texttt{POST /data}):} Recibe el JSON enviado por Lazarus, calcula la fecha y hora exacta (\texttt{DateTime}) y guarda el registro en la tabla \texttt{mediciones} del archivo \texttt{datos\_ambientales.db}.
        \item \textbf{API REST en formato JSON (\texttt{GET /data}, \texttt{/data/ultimo}, \texttt{/data/uma/<id>}):} Permite consultar desde el navegador los registros guardados en SQLite directamente en formato JSON.
        \item \textbf{Gráfica Matplotlib (\texttt{GET /grafico.png} y \texttt{GET /}):} Consulta las últimas 20 mediciones en SQLite y genera con \textbf{Matplotlib} una figura con \textbf{3 Subplots} (evolución de MP, histograma de 6 diámetros y evolución de Temperatura/Humedad), mostrándola en el panel web (\texttt{templates/index.html}).
    \end{itemize}
\end{enumerate}

\section{Contrato de Datos: La Trama de 12 Valores}

Tanto el simulador serial, como Lazarus, el JSON de envío y la tabla de SQLite manejan exactamente los mismos \textbf{12 datos enteros} en el mismo orden.

Por ejemplo, ante la trama:
\begin{center}
\texttt{4, 62, 48, 71, 5593, 1398, 9349, 1868, 5142, 6684, 18, 45}
\end{center}

\begin{table}[h!]
\centering
\small
\begin{tabular}{@{}clllc@{}}
\toprule
\textbf{Posición} & \textbf{Nombre en Código / JSON} & \textbf{Qué representa físicamente} & \textbf{Dónde se muestra en Lazarus} & \textbf{Rango} \\ \midrule
\texttt{0} & \texttt{id\_uma} & Estación Ambiental (UMA 1 a 5) & Filtro \texttt{RbUma1..RbUma5} & $1 \dots 5$ \\
\texttt{1} & \texttt{mp1\_0} & Material Particulado PM 1.0 ($\mu\text{g/m}^3$) & Curva azul (\texttt{SerieMP1}) & $5 \dots 100$ \\
\texttt{2} & \texttt{mp2\_5} & Material Particulado PM 2.5 ($\mu\text{g/m}^3$) & Curva verde (\texttt{SerieMP25}) & $5 \dots 100$ \\
\texttt{3} & \texttt{mp10} & Material Particulado PM 10 ($\mu\text{g/m}^3$) & Curva roja (\texttt{SerieMP10}) & $5 \dots 100$ \\
\texttt{4} & \texttt{part\_03} & Conteo partículas de $0.3\,\mu\text{m}$ & Barra 0 del Histograma & $10 \dots 10000$ \\
\texttt{5} & \texttt{part\_05} & Conteo partículas de $0.5\,\mu\text{m}$ & Barra 1 del Histograma & $10 \dots 10000$ \\
\texttt{6} & \texttt{part\_10} & Conteo partículas de $1.0\,\mu\text{m}$ & Barra 2 del Histograma & $10 \dots 10000$ \\
\texttt{7} & \texttt{part\_25} & Conteo partículas de $2.5\,\mu\text{m}$ & Barra 3 del Histograma & $10 \dots 10000$ \\
\texttt{8} & \texttt{part\_50} & Conteo partículas de $5.0\,\mu\text{m}$ & Barra 4 del Histograma & $10 \dots 10000$ \\
\texttt{9} & \texttt{part\_100} & Conteo partículas de $10.0\,\mu\text{m}$ & Barra 5 del Histograma & $10 \dots 10000$ \\
\texttt{10} & \texttt{temperatura} & Temperatura ambiente ($^\circ\text{C}$) & Cuadro \texttt{LblTemperatura} & $5 \dots 30$ \\
\texttt{11} & \texttt{humedad} & Humedad relativa del aire ($\%$) & Cuadro \texttt{LblHumedad} & $20 \dots 70$ \\ \bottomrule
\end{tabular}
\caption{Estructura de los 12 valores de cada medición.}
\end{table}

\section{El ``Por qué'' de las Cosas (Explicación Conceptual por Módulo)}

\subsection{1. Simulador Serial (\texttt{2\_simulador\_serial/simulador\_serial.py})}

\begin{itemize}[leftmargin=1.5em]
    \item \textbf{¿Por qué se creó la función \texttt{variar\_valor} en vez de usar \texttt{randint} puro?}
    Si en cada lectura se generara un número 100\% al azar entre $5$ y $30^\circ\text{C}$, la temperatura saltaría de $6^\circ\text{C}$ a $29^\circ\text{C}$ en medio segundo, lo cual se ve irreal. Guardando un diccionario \texttt{estado\_estaciones} para las 5 UMAs y sumando o restando una pequeña variación en cada paso (sin salirse nunca de los límites de la pauta), las curvas se mueven de forma suave como un sensor real.
    
    \item \textbf{¿Por qué al enviar por serial se escribe \texttt{(trama + "\textbackslash r\textbackslash n").encode("utf-8")}?}
    Por dos motivos:
    \begin{enumerate}
        \item En Python 3 la librería \texttt{pyserial} transmite bytes crudos por el puerto COM, no cadenas \texttt{str}, por lo que \texttt{.encode("utf-8")} transforma el texto a bytes.
        \item El salto de línea \texttt{\textbackslash r\textbackslash n} al final es fundamental porque en Lazarus el componente \texttt{TLazSerial} tiene activada la propiedad \texttt{RcvLineCRLF = True}, la cual espera hasta encontrar un salto de línea para entregar la trama completa sin cortarla a la mitad.
    \end{enumerate}
\end{itemize}

\subsection{2. Interfaz Gráfica en Lazarus (\texttt{3\_lazarus\_datalogger/main.pas})}

\begin{itemize}[leftmargin=1.5em]
    \item \textbf{¿Cómo se lee y desarma la trama serial en \texttt{PuertoSerialRxData}?}
    Cuando llega información al puerto \texttt{COM2}, se dispara automáticamente el evento \texttt{PuertoSerialRxData}. En vez de cortar el texto a mano buscando posiciones de comas, usamos un objeto \texttt{TStringList} (\texttt{elementos\_linea}) configurando \texttt{Delimiter := ','}. Eso separa automáticamente el texto en una lista de elementos. Luego verificamos que tenga exactamente 12 elementos (\texttt{Count = 12}) y los convertimos a entero de forma segura con \texttt{TryStrToInt}.
    
    \item \textbf{¿Cómo funciona el efecto \emph{Scrolling Left} (desplazamiento a la izquierda)?}
    El gráfico de líneas muestra una ventana de 11 muestras en el eje horizontal (posiciones \texttt{0} a \texttt{10}):
    \begin{enumerate}
        \item Cuando llega un dato nuevo de una UMA, un ciclo \texttt{for j := 0 to 9} mueve cada valor antiguo un puesto hacia la izquierda: \texttt{historial\_mp1[uma, j] := historial\_mp1[uma, j + 1]}. El dato más viejo (posición \texttt{0}) sale del gráfico.
        \item El dato recién llegado se guarda en la última posición de la derecha (\texttt{historial\_mp1[uma, 10]}).
        \item Finalmente, \texttt{DibujarEstacionActual} limpia las series y vuelve a pintar los puntos del \texttt{0} al \texttt{10}, logrando que la curva se desplace continuamente de derecha a izquierda.
    \end{enumerate}
    
    \item \textbf{¿Por qué los arreglos de historial tienen dos dimensiones \texttt{[1..5, 0..10]}?}
    La primera dimensión (\texttt{1..5}) representa el número de estación UMA y la segunda (\texttt{0..10}) las 11 muestras en el tiempo. Así, aunque en pantalla estemos mirando la \texttt{UMA 1}, el programa sigue guardando en segundo plano los datos que llegan de las \texttt{UMA 2, 3, 4 y 5}. Cuando el usuario cambia el \texttt{RadioButton} a otra UMA, los gráficos no parten de cero ni se borran, sino que muestran de inmediato el historial acumulado de esa estación.
    
    \item \textbf{¿Cómo funcionan los RadioButtons \texttt{MP 01}, \texttt{MP 25} y \texttt{MP 10}?}
    Los tres gráficos de línea están dibujados al mismo tiempo, pero si mostráramos los carteles con números (\texttt{Marks}) en las 3 curvas a la vez, se taparían entre sí. El procedimiento \texttt{ActualizarEtiquetasMP} activa \texttt{Marks.Visible := True} únicamente en la curva cuyo \texttt{RadioButton} está marcado.
    
    \item \textbf{¿Por qué se usa un \texttt{TBitmap} en caché (\texttt{bitmap\_fondo\_cache}) para el fondo del gráfico?}
    El gráfico \texttt{Datos Ambientales} dibuja un paisaje de fondo mediante el evento \texttt{OnAfterDrawBackWall} (\texttt{DibujarFondoGrafico}). Si re-escaláramos el archivo \texttt{fondo\_grafico.jpg} cada vez que entra un dato serial, el procesador se saturaría y la ventana se sentiría lenta. Por eso, la primera vez que se dibuja (o solo si cambia el tamaño de la ventana), se escala una sola vez dentro de \texttt{bitmap\_fondo\_cache}, y en las lecturas siguientes se hace una copia directa desde memoria RAM (\texttt{ACanvas.Draw}), que toma menos de $0.1\,\text{ms}$.
    
    \item \textbf{¿Cómo envía Lazarus los datos al servidor Flask (\texttt{EnviarDatosFlask})?}
    Usa las unidades nativas de Free Pascal \texttt{fpjson} y \texttt{fphttpclient}:
    \begin{enumerate}
        \item Crea un \texttt{TJSONObject} y agrega los 12 pares clave-valor (\texttt{'id\_uma'}, \texttt{'mp1\_0'}, etc.).
        \item Crea un cliente \texttt{TFPHTTPClient}, agrega el encabezado \texttt{Content-Type: application/json} (indispensable para que Flask entienda que el cuerpo de la petición es un JSON) y envía el texto mediante \texttt{Post('http://127.0.0.1:7070/data')}.
        \item Todo el envío está envuelto en un bloque \texttt{try..except}: si el servidor Flask está apagado o se reinicia, Lazarus captura el error sin congelarse ni cerrarse y solo avisa en la barra de estado inferior.
    \end{enumerate}
\end{itemize}

\subsection{3. Servidor Flask, SQLite y Matplotlib (\texttt{1\_servidor\_flask/servidor.py})}

\begin{itemize}[leftmargin=1.5em]
    \item \textbf{¿Cómo se guarda el registro con \texttt{DateTime} en SQLite?}
    Al arrancar el servidor, \texttt{inicializar\_bd()} crea automáticamente el archivo \texttt{datos\_ambientales.db} y la tabla \texttt{mediciones}. Cada vez que llega una petición \texttt{POST} a \texttt{/data}, el servidor toma la hora exacta del sistema con \texttt{datetime.now().strftime("\%Y-\%m-\%d \%H:\%M:\%S")} y ejecuta un \texttt{INSERT INTO mediciones} guardando la fecha/hora junto con los 12 valores ambientales.
    
    \item \textbf{¿Por qué se ejecuta \texttt{PRAGMA journal\_mode=WAL;} al abrir SQLite?}
    \texttt{WAL} (\emph{Write-Ahead Logging}) permite que SQLite reciba escrituras desde Lazarus (\texttt{POST /data}) al mismo tiempo que el navegador realiza lecturas (\texttt{GET /data} o \texttt{GET /grafico.png}) sin bloquear la base de datos.
    
    \item \textbf{¿Por qué se usa \texttt{conexion.row\_factory = sqlite3.Row}?}
    Por defecto, SQLite en Python devuelve tuplas sin nombre \texttt{(1, '2026-11-06...', 45, ...)}. Al usar \texttt{sqlite3.Row}, cada fila se puede convertir directamente en un diccionario con \texttt{dict(fila)} y transformarse en respuesta JSON con \texttt{jsonify(...)}.
    
    \item \textbf{¿Cómo se genera el gráfico de Matplotlib de 3 Subplots (\texttt{generar\_grafico})?}
    \begin{enumerate}
        \item Al inicio del archivo se declara \texttt{matplotlib.use('Agg')}. Esto activa el motor sin interfaz gráfica de Matplotlib para que genere archivos PNG en el servidor sin intentar abrir ventanas emergentes de escritorio.
        \item Se consultan en SQLite las últimas 20 mediciones (\texttt{ORDER BY id DESC LIMIT 20}) y se invierte la lista (\texttt{registros.reverse()}) para que en el eje horizontal el tiempo avance de más antiguo (izquierda) a más reciente (derecha).
        \item Se divide la figura en 3 filas (\texttt{plt.subplot(3, 1, 1)}, \texttt{(3, 1, 2)} y \texttt{(3, 1, 3)}):
        \begin{itemize}
            \item \textbf{Subplot 1:} Grafica con \texttt{plt.plot} las 3 curvas históricas de \texttt{MP 1.0}, \texttt{MP 2.5} y \texttt{MP 10}.
            \item \textbf{Subplot 2:} Toma la medición más reciente (\texttt{registros[-1]}) y dibuja con \texttt{plt.bar} el histograma de 6 barras de colores para los diámetros $0.3\,\mu\text{m}$ a $10.0\,\mu\text{m}$, colocando el valor sobre cada barra con \texttt{plt.text}.
            \item \textbf{Subplot 3:} Grafica con \texttt{plt.plot} la evolución histórica de Temperatura ($^\circ\text{C}$) y Humedad Relativa ($\%$).
        \end{itemize}
        \item Guarda la imagen en \texttt{1\_servidor\_flask/static/grafico.png}, cierra la figura con \texttt{plt.close(fig)} para liberar memoria RAM y la ruta \texttt{/grafico.png} la entrega al navegador con \texttt{send\_file}.
    \end{enumerate}
\end{itemize}

\newpage
\section{Mapa Rápido para la Defensa: ¿Dónde está cada cosa?}

Si durante la presentación el profesor te pide mostrar una parte específica del código, usa esta tabla para ir directo al archivo y función correspondiente:

\begin{table}[h!]
\centering
\small
\begin{tabular}{@{}p{5.8cm}p{4.6cm}p{5.2cm}@{}}
\toprule
\textbf{Si el profesor pregunta o pide ver...} & \textbf{Archivo} & \textbf{Función / Procedimiento y Líneas} \\ \midrule
1. ¿Dónde se genera la trama aleatoria de 12 datos en Python? & \texttt{2\_simulador\_serial/\newline simulador\_serial.py} & Función \texttt{generar\_trama(id\_uma)} (líneas 35--60). \\ \addlinespace
2. ¿Dónde se abre el puerto \texttt{COM2} y se envía por serial en Python? & \texttt{2\_simulador\_serial/\newline simulador\_serial.py} & Bloque \texttt{if \_\_name\_\_ == "\_\_main\_\_":} (líneas 63--88). \\ \addlinespace
3. ¿Dónde se conecta o desconecta el puerto \texttt{COM2} en Lazarus? & \texttt{3\_lazarus\_datalogger/\newline main.pas} & Procedimiento \texttt{ChkConectarSerialChange} (líneas 241--267). \\ \addlinespace
4. ¿Dónde recibe Lazarus la trama del puerto serial y separa las comas? & \texttt{3\_lazarus\_datalogger/\newline main.pas} & Procedimiento \texttt{PuertoSerialRxData} (líneas 380--430). \\ \addlinespace
5. ¿Dónde se hace el \emph{Scrolling Left} (correr los datos a la izquierda)? & \texttt{3\_lazarus\_datalogger/\newline main.pas} & Dentro de \texttt{PuertoSerialRxData} (líneas 447--458) y dibujo en \texttt{DibujarEstacionActual} (líneas 312--335). \\ \addlinespace
6. ¿Dónde se activan/desactivan los carteles numéricos \texttt{MP 01/25/10}? & \texttt{3\_lazarus\_datalogger/\newline main.pas} & Procedimientos \texttt{ActualizarEtiquetasMP} y \texttt{CambiarEtiquetasMP} (líneas 279--295). \\ \addlinespace
7. ¿Dónde arma Lazarus el JSON y hace el \texttt{POST} con \texttt{TFPHTTPClient}? & \texttt{3\_lazarus\_datalogger/\newline main.pas} & Procedimiento \texttt{EnviarDatosFlask} (líneas 345--377). \\ \addlinespace
8. ¿Dónde se guardan las capturas PNG de los gráficos de Lazarus? & \texttt{3\_lazarus\_datalogger/\newline main.pas} & \texttt{BtnGuardarSeriesClick} y \texttt{BtnGuardarHistoClick} (líneas 485--517). \\ \addlinespace
9. ¿Dónde se crea la base de datos SQLite y la tabla \texttt{mediciones}? & \texttt{1\_servidor\_flask/\newline servidor.py} & Función \texttt{inicializar\_bd()} (líneas 28--52). \\ \addlinespace
10. ¿Dónde recibe Flask el JSON, le agrega \texttt{DateTime} y hace el \texttt{INSERT}? & \texttt{1\_servidor\_flask/\newline servidor.py} & Función \texttt{manejar\_data()} en el bloque \texttt{if request.method == 'POST':} (líneas 151--196). \\ \addlinespace
11. ¿Dónde están los endpoints de la API REST que devuelven JSON al navegador? & \texttt{1\_servidor\_flask/\newline servidor.py} & \texttt{GET /data} (líneas 198--221),\newline \texttt{GET /data/ultimo} (líneas 224--247),\newline \texttt{GET /data/uma/<id>} (líneas 250--264). \\ \addlinespace
12. ¿Dónde se crea el gráfico Matplotlib de 3 Subplots? & \texttt{1\_servidor\_flask/\newline servidor.py} & Función \texttt{generar\_grafico()} (líneas 55--148) y ruta \texttt{/grafico.png} (líneas 267--272). \\ \bottomrule
\end{tabular}
\caption{Guía rápida de ubicación en el código para preguntas de la defensa.}
\end{table}

\section{Orden Exacto para Iniciar el Proyecto en la Defensa}

\begin{enumerate}[leftmargin=1.5em]
    \item \textbf{Abrir VSPE:} Crear (o cargar) el dispositivo \textbf{\texttt{Connector}} en el puerto \textbf{\texttt{COM2}} y darle al botón \textbf{Play} (\emph{Start Emulation}).
    \item \textbf{Abrir Lazarus IDE:} Abrir \texttt{3\_lazarus\_datalogger/datalogger.lpi}, presionar \textbf{F9} para ejecutar la ventana, y marcar arriba las casillas \textbf{\texttt{COM On/Off}} y \textbf{\texttt{Envía a Flask}}.
    \item \textbf{Terminal 1 (Servidor Flask):}
    \begin{verbatim}
    cd 1_servidor_flask
    python servidor.py
    \end{verbatim}
    \item \textbf{Terminal 2 (Simulador Serial):}
    \begin{verbatim}
    cd 2_simulador_serial
    python simulador_serial.py
    \end{verbatim}
    \item \textbf{Navegador Web:}
    \begin{itemize}
        \item Panel principal y gráfica Matplotlib (3 Subplots): \url{http://127.0.0.1:7070/}
        \item Consultas JSON en el navegador: \url{http://127.0.0.1:7070/data}, \url{http://127.0.0.1:7070/data/ultimo} y \url{http://127.0.0.1:7070/data/uma/1}
    \end{itemize}
\end{enumerate}

\end{document}
